% Generated draft: requires mathematical and source review.
\documentclass[11pt]{article}
\usepackage[margin=25mm]{geometry}
\usepackage{fontspec}
\setmainfont{Latin Modern Roman}
\usepackage{amsmath,amssymb,mathtools,unicode-math}
\setmathfont{Latin Modern Math}
\usepackage{enumitem,needspace,xurl,hyperref}
\setlength{\parindent}{0pt}
\setlength{\parskip}{5pt}
\setlength{\emergencystretch}{3em}
\newcommand{\N}{\mathbb{N}}
\newcommand{\Z}{\mathbb{Z}}
\newcommand{\Q}{\mathbb{Q}}
\newcommand{\R}{\mathbb{R}}
\newcommand{\C}{\mathbb{C}}
\newcommand{\F}{\mathbb{F}}
\newcommand{\A}{\mathbb{A}}
\newcommand{\PP}{\mathbb{P}}
\newcommand{\E}{\mathbb{E}}
\newcommand{\eps}{\varepsilon}
\newcommand{\dd}{\,\mathrm{d}}
\begin{document}
\section*{UnsolvedMath ID: 30001160}
\textbf{Source problem code:} catalogue entry
\subsection{Definitions and mathematical statement}
An eigenvalue isoperimetric problem is understood here only schematically as an optimization problem for a quantity determined by eigenvalues, subject to geometric constraints. For a specified admissible class \(\mathcal A\) and a specified objective \(J\colon\mathcal A\to\R\cup\{+\infty\}\), an object \(X_*\in\mathcal A\) is a global minimizer if
\[
J(X_*)=\inf_{X\in\mathcal A}J(X).
\]
This notation describes the minimizer terminology only; the supplied record does not specify \(\mathcal A\), \(J\), or the underlying eigenvalue problem.

The source question is: \emph{Can an isoperimetric problem for eigenvalues have a fractal minimizer or other fractal solution?}

The question asks whether such a phenomenon can occur, without specifying a particular optimization problem. A precise instance would require an ambient space, an operator and its realization, boundary conditions where applicable, an eigenvalue objective, geometric constraints, an admissible class, and a definition of fractality. In particular, the record does not say whether fractality concerns the optimizing object itself, its boundary, or another associated set. The phrase ``other fractal solution'' also has no specified mathematical meaning. Consequently, the displayed minimization formula is only a schematic description of one branch of the question; a fully quantified mathematical statement cannot be recovered from the supplied record without additional choices.
\subsection{Short English statement}
Can an eigenvalue optimization problem subject to geometric constraints admit a fractal minimizer, or another kind of fractal solution? The source leaves the operator, constraints, admissible objects, meaning of fractality, and alternative notion of solution unspecified.
\subsection{Sources}
\begin{itemize}
\item Supplied local UnsolvedMath registry record, permanent ID \(30001160\), title ``Fractal Optimizers in Spectral Isoperimetric Problems,'' exact source problem code \texttt{catalogue entry}. The record belongs to the Oberwolfach Reports Open Problems collection and gives the proposed year as 2009. No individual proposer is supplied.
\item \emph{Low Eigenvalues of Laplace and Schrödinger Operators} (2009), Oberwolfach Reports, pp.~355--428, DOI 10.4171/owr/2009/06. Supplied source-report and literature-source link: \url{https://doi.org/10.4171/owr/2009/06}. The collection describes these reports as documenting mathematical workshops and their open-problem discussions.
\item Supplied literature source labeled ``regularity of optimal spectral domains'': \url{https://doi.org/10.1007/s00205-014-0801-6}. No authors, title beyond this descriptive label, or relevant theorem are supplied in the record.
\item Literature assessment dated \texttt{2026-08-21}: the supplied curation labels the problem open and describes existence of a fractal eigenvalue optimizer as an unresolved regularity and attainment question. It states that later spectral-shape theory provides regularity for important relaxed problems without generally resolving this question. This is supplied dated curation, not an independently verified literature assessment; its time-sensitive status and cited claims require independent verification.
\item The registry marks the statement as ``exact'' and reports that its mathematical content was checked against the official Oberwolfach report, with notation, context, and question form normalized. This is a supplied source-verification assertion, not verification performed for this draft.
\end{itemize}
\end{document}
